\documentclass[a4paper,10pt]{article}

% Packages
\usepackage[bottom=1.5in, top=1in, left=1in, right=1in]{geometry} % set smaller margins (default ~1.5in)
\usepackage{fancyhdr}   % for header/footer customization
\usepackage{graphicx}   % for including images
\usepackage{fontawesome5}
\usepackage{enumitem}
\usepackage{multicol}

\setlist{itemsep=0pt, topsep=0pt}

% Header setup
\pagestyle{fancy}
\fancyhf{} % clear all header and footer fields

% Insert logos: left and right
\fancyhead[L]{\includegraphics[width=2.8cm]{images/upc-logo.png}}
\fancyhead[R]{\includegraphics[width=2.8cm]{images/dcs-logo-full.png}}

% Center text in header (two lines: top small, bottom main title)
\fancyhead[C]{%
  \textbf{\Large University of the Philippines Cebu} \\[4pt]%
  \small College of Science \\%
  \textbf{\small Department of Computer Science} \\%
  \faIcon{phone} 032 232 8187 \faIcon{envelope} dcs.upcebu@up.edu.ph

}

% Keep horizontal line below header
\renewcommand{\headrulewidth}{0.4pt}

% Increase top margin so body text starts clearly below the line
\setlength{\headheight}{85pt}  % must be >= height of logos/text
\setlength{\footskip}{25pt}    % distance from bottom of text to footer
\setlength{\headsep}{15pt}     % distance from header to text

% Footer
\fancyfoot[C]{\thepage} % page number at center bottom
\renewcommand{\footrulewidth}{0.4pt} % line above footer

\begin{document}


\begin{center}
    \textbf{\Large COURSE GUIDE}
\end{center}
\vspace{0.5em}

\begin{tabular}{ll}
\textbf{Course Number}   & CMSC 178IP \\
\textbf{Course Title}    & Digital Image Processing \\
\textbf{Number of Units} & 3 units (3 hours lecture) \\
\textbf{Instructor}      & Noel Jeffrey Pinton \\
\end{tabular}

\vspace{1.5em}

\begin{center}
    \textbf{\Large COURSE OUTLINE}
\end{center}

\begin{enumerate}[label=\Roman*.]
    \item \textbf{Introduction}
    \begin{enumerate}[label=\alph*.]
        \item Overview and Applications
        \item Digital Image Fundamentals
        \begin{enumerate}[label=\roman*.]
            \item Image as a 2D function $f(x, y)$
            \item Pixels, sampling, and quantization
            \item Image types and storage formats
        \end{enumerate}
    \end{enumerate}

    \item \textbf{Image Formation, Representation, and Compression}
    \begin{enumerate}[label=\alph*.]
        \item Image Coding and Compression
        \begin{enumerate}[label=\roman*.]
            \item Lossless compression: Huffman, Run-length, LZW
            \item Lossy compression: JPEG, JPEG2000, WebP
        \end{enumerate}
        \item Imaging Sensors and Systems
        \begin{enumerate}[label=\roman*.]
            \item CCD, CMOS, PET, CT
            \item Noise sources in image acquisition
        \end{enumerate}
    \end{enumerate}

    \item \textbf{Image Enhancement}
    \begin{enumerate}[label=\alph*.]
        \item Spatial Domain Methods
        \begin{enumerate}[label=\roman*.]
            \item Spatial filtering: smoothing, sharpening
            \item Histogram equalization and specification
        \end{enumerate}
        \item Frequency Domain Methods
        \begin{enumerate}[label=\roman*.]
            \item 2D Fourier Transform, frequency representation
            \item Low-pass, high-pass, band-pass filtering
        \end{enumerate}
        \item Noise Reduction Techniques
        \begin{enumerate}[label=\roman*.]
            \item Linear filtering (mean, Gaussian)
            \item Nonlinear filtering (median, bilateral)
        \end{enumerate}
    \end{enumerate}

    \item \textbf{Image Restoration and Geometric Processing}
    \begin{enumerate}[label=\alph*.]
        \item Image Filtering and Convolution
        \begin{enumerate}[label=\roman*.]
            \item Convolution vs. correlation
            \item Spatial kernels: blurring, sharpening, edge enhancement
        \end{enumerate}
        \item Geometric Transformations
        \begin{enumerate}[label=\roman*.]
            \item 2D Scaling, rotation, translation, shearing
            \item Affine and projective transformations
            \item Image registration
        \end{enumerate}
    \end{enumerate}

    \item \textbf{Feature Extraction and Representation}
    \begin{enumerate}[label=\alph*.]
        \item Edge and Line Detection
        \begin{enumerate}[label=\roman*.]
            \item Gradient operators (Sobel, Prewitt, Roberts)
            \item Hough transform for lines and circles
        \end{enumerate}
        \item Texture and Local Features
    \end{enumerate}

    \item \textbf{Image Segmentation and Morphology}
    \begin{enumerate}[label=\alph*.]
        \item Segmentation Methods
        \begin{enumerate}[label=\roman*.]
            \item Thresholding (global, local, Otsu's method)
            \item Region-based segmentation (region growing, watershed)
        \end{enumerate}
        \item Morphological Processing
        \begin{enumerate}[label=\roman*.]
            \item Binary morphology: erosion, dilation, opening, closing
            \item Morphology for grayscale images
        \end{enumerate}
    \end{enumerate}

    \item \textbf{Computer Vision and Deep Learning Approaches}
    \begin{enumerate}[label=\alph*.]
        \item Convolutional Neural Networks (CNNs)
        \begin{enumerate}[label=\roman*.]
            \item CNN architecture basics: convolution, pooling, fully connected layers
            \item Activation functions (ReLU, sigmoid, tanh)
            \item Pooling and subsampling
        \end{enumerate}
        \item Applications
        \begin{enumerate}[label=\roman*.]
            \item Image classification (MNIST, CIFAR-10)
            \item Image segmentation
            \item Object detection (YOLO, Faster R-CNN)
        \end{enumerate}
        \item Generative Models
        \begin{enumerate}[label=\roman*.]
            \item Variational Autoencoders (VAE)
            \item Generative Adversarial Networks (GAN)
        \end{enumerate}
    \end{enumerate}
\end{enumerate}

\vspace{1.5em}

\newpage

\section*{Group Project Guidelines -- CMSC 178IP (Digital Image Processing)}

\subsection*{1. Group Formation}
\begin{itemize}
    \item Students should organize themselves into \textbf{groups of up to 4 members}.
    \item Smaller groups (2--3) are allowed, but workload expectations remain the same regardless of group size.
    \item Each group must register their members with the instructor by \textbf{Midterms week}.
    \item Groups are encouraged to form with diverse skills (programming, mathematics, presentation skills, documentation).
\end{itemize}

\subsection*{2. Project Proposal}
\textbf{Due by Midterms.}

\noindent The proposal should include:
\begin{enumerate}
    \item Title and objective of the project.
    \item Motivation: why is this problem interesting/important?
    \item Background: relevant references (papers, applications, prior work).
    \item Dataset to be used (must be accessible and ethical to use).
    \item Methodology: which image processing or computer vision techniques will be applied.
    \item Expected outcomes and possible applications.
\end{enumerate}

\subsection*{3. Project Development Process}

\subsubsection*{1. Data Gathering and Preparation}
Collect suitable image data (public datasets preferred, or self-collected with consent). Perform preprocessing (resizing, normalization, annotation if needed).

\subsubsection*{2. Method Implementation}
Implement at least two non-trivial techniques from the course (e.g., morphological processing, edge detection, CNNs). Justify parameter choices (filters, thresholds, architectures).

\subsubsection*{3. Evaluation and Analysis}
Define evaluation metrics (PSNR, SSIM, accuracy, F1-score, IoU, etc.). Critically analyze results, including strengths and limitations. Compare against baseline methods where possible.

\subsection*{4. Deliverables}
\begin{enumerate}
    \item \textbf{Project Report} with abstract, introduction, related work, methodology, results, discussion, conclusion, and references.
    \item \textbf{Code Repository} (Python preferred, with README and requirements) that reproduces results.
    \item \textbf{Oral Presentation} (15--20 minutes during TBD). Each member must speak.
\end{enumerate}

\subsection*{5. Grading Breakdown}
\begin{itemize}
    \item Proposal: 10\%
    \item Implementation \& Technical Depth: 30\%
    \item Results \& Analysis: 20\%
    \item Report (clarity, structure, rigor): 20\%
    \item Presentation \& Delivery: 20\%
\end{itemize}

\subsection*{6. Expectations}
\begin{itemize}
    \item Projects should go beyond toy examples. Examples:
    \begin{itemize}
        \item Image denoising with classical filters vs. CNN-based denoising.
        \item Medical image segmentation using thresholding vs. morphological operations.
        \item Object classification using handcrafted features vs. CNN.
    \end{itemize}
    \item Creativity and originality are highly encouraged.
    \item Academic integrity must be maintained: no plagiarism, proper citation, and acknowledgment of external code or datasets.
\end{itemize}

\end{document}
